\documentclass{article}
\usepackage[UTF8]{ctex}
\usepackage{amsmath,amssymb,amsfonts}
\usepackage{enumitem}
\usepackage{geometry}
\geometry{a4paper, margin=1in}
\usepackage{graphicx}
% === 错题本排版优化（不参与题目标识） ===
\setlength{\parindent}{0pt}
\setlength{\parskip}{0.35em plus 0.10em minus 0.08em}
\setlength{\emergencystretch}{3em}
\setlist{itemsep=0.25em, topsep=0.40em, parsep=0pt}
\allowdisplaybreaks[2]
% === 错题本排版优化结束 ===
\begin{document}
\section*{山东大学2025年高等代数真题}
\section*{114．山东大学2025年线性代数与常微分方程试卷}
线性代数\\

\subsection{题目1}
1．（15 分）设齐次线性方程组

\[
\left\{\begin{array}{c}
a_{11} x_{1}+a_{12} x_{2}+\cdots+a_{1 n} x_{n}=0 \\
a_{21} x_{1}+a_{22} x_{2}+\cdots+a_{2 n} x_{n}=0 \\
\cdots \cdots \\
a_{n-1,1} x_{1}+a_{n-1,2} x_{2}+\cdots+a_{n-1, n} x_{n}=0
\end{array}\right.
\]

证明：\(X_{0}=\left(x_{1_{0}}, x_{2_{0}}, \cdots, x_{n_{0}}\right)^{\mathrm{T}}\) 为此方程组的解，其中

\[
\begin{gathered}
x_{1_{0}}=\left|\begin{array}{cccc}
a_{12} & a_{13} & \cdots & a_{1 n} \\
a_{22} & a_{23} & \cdots & a_{2 n} \\
\vdots & \vdots & & \vdots \\
a_{n-1,2} & a_{n-1,3} & \cdots & a_{n-1, n}
\end{array}\right|, x_{2_{0}}=-\left|\begin{array}{cccc}
a_{11} & a_{13} & \cdots & a_{1 n} \\
a_{21} & a_{23} & \cdots & a_{2 n} \\
\vdots & \vdots & & \vdots \\
a_{n-1,1} & a_{n-1,3} & \cdots & a_{n-1, n}
\end{array}\right|, \\
\cdots, x_{n_{0}}=(-1)^{n+1}\left|\begin{array}{cccc}
a_{11} & a_{12} & \cdots & a_{1, n-1} \\
a_{21} & a_{22} & \cdots & a_{2, n-1} \\
\vdots & \vdots & & \vdots \\
a_{n-1,1} & a_{n-1,2} & \cdots & a_{n-1, n-1}
\end{array}\right| .
\end{gathered}
\]

且若 \(X_{0} \neq 0\) ，则方程组的任一解可表示为 \(k X_{0}\) ，其中 \(k\) 为常数。\\

\subsection{题目2}
2．（15 分）设 \(A\) 为 \(n\) 阶可逆反对称矩阵，\(b\) 为 \(n\) 维列向量，\(B=\left(\begin{array}{cc}A & b \\ b^{\mathrm{T}} & 0\end{array}\right)\) 。证明：\\
（1）\(n\) 为偶数．\\
（2）矩阵 \(B\) 的秩 \(r(B)=n\) 。\\

\subsection{题目3}
3．（10 分）设 \(A\) 为 \(n\) 阶正定矩阵，\(B\) 为 \(n \times m\) 阶实矩阵，证明：若 \(r(B)=m\) ，则 \(B^{\mathrm{T}} A B\) 为正定矩阵。常微岔夯程正明如下（1）和（2）是等价的。\\
（1）（文除）\(A\) 求是加程矩阵．\(\left.x y^{2}\right) \mathrm{d} x+\left(x^{2} y+y^{3}\right) \mathrm{d} y=0\) 的通解。\\
\textit{[image omitted]}解 \(\mid y=\psi(\) 时），在整每实输死素等再该元菊的代数余子式乘以 -1 。\\

\subsection{题目4}
3．（1＆分））投解下証方，程 者具有相司特征多项式的） 3 阶复方阵，证明：\(A, B, C, D\) 中必有两个矩阵相似。（10 分）解方程 \(y^{\prime \prime}-2 y^{\prime}+2 y=x e^{x} \sin x\) 。\\

\subsection{题目5}
6．（20 分）设结定戒数域程组將 线性空间其中笧数零除空间，中船无素十的线性是㱩卖周期函数，\\
\textit{[image omitted]}使）䳌im \((\sigma(W))+\operatorname{dim}(\operatorname{Ker} \sigma \cap W)=\operatorname{dim} W\) ．\\
（2） \(\operatorname{dim} W \leq \operatorname{dim}\left(\sigma^{-1}(W)\right) \leq \operatorname{dim} W+\operatorname{dim}(\operatorname{Ker} \sigma)\) ．

\[
\mathbf{X}(t+T)=\mathbf{X}(t) \mathbf{B}
\]
\end{document}
